%% 5-15 sin・cosを微分すると？（第11回：ページ制限なし。説明・図・例題を書き足した）
\newcommand\bkvomidashi[1]{\par\Needspace{6\baselineskip}\noindent\textbf{#1}\par}
\bkchapter{5}{15}{$\sin$・$\cos$を微分すると？}{\bklevel{4}}{sin・cosを微分すると？}
\bkmichishirube{4-2,4-8,5-6}{5-17,6-6,6-10}

\begin{goal}
\begin{itemize}[label={\color{bkTheme}●}, leftmargin=1.6\zw, itemsep=0pt, topsep=1pt]
\item \bklevel{3}\enspace $(\sin t)'=\cos t$、$(\cos t)'=-\sin t$ を、円運動で説明できる。
\item \bklevel{4}\enspace $\sin\omega t$、$\cos\omega t$ を微分・積分できる。
\item \bklevel{4}\enspace 単振動の速度・加速度を求め、$\ddot{x}=-\omega^2x$ を確かめられる。
\item \bklevel{4}\enspace $\sin^2$ の平均から、交流の実効値を求められる。\bkidx{こうりゅう}{交流}
\end{itemize}
\end{goal}

\begin{hatena}
\noindent\begin{minipage}[c]{0.58\linewidth}
観覧車のゴンドラは、一定の速さで円をえがいて回ります。でも、地面から見た高さの変わり方は一定ではありません。いちばん下や上ではほとんど変わらず、横を通るときに、いちばん速く上がったり下がったりします。

高さは $\sin$ や $\cos$ で表せました（4-8）。その「変わる速さ」は、どんな式になるのでしょう。
\end{minipage}\hfill
\begin{minipage}[c]{0.39\linewidth}
\centering
\begin{tikzpicture}[x=0.78mm, y=0.78mm]
  \draw[bkGray, line width=1pt] (-20,-22) -- (20,-22);
  \draw[bkInk!60, line width=0.8pt] (0,0) circle (17);
  \draw[bkInk!60, line width=0.8pt] (0,0) -- (-8,-22) (0,0) -- (8,-22);
  \foreach \a in {0,45,...,315} {\draw[bkInk!40, line width=0.5pt] (0,0) -- (\a:17); \fill[bkRoute2] (\a:17) ++(-1.8,-2.6) rectangle ++(3.6,3);}
  \fill[bkOrange] (0:17) ++(-1.8,-2.6) rectangle ++(3.6,3);
  \draw[bk arrow, bkOrange, line width=1pt] (21,-3) -- (21,7);
  \node[font=\sffamily\scriptsize, anchor=west, align=left] at (22,2) {ここで\\いちばん\\速く上がる};
\end{tikzpicture}
\end{minipage}
\end{hatena}

$\sin$ と $\cos$ は、回るもの・ゆれるもの・波を表す関数です（4-8）。この章では、これらを微分すると、また $\sin$ と $\cos$ が出てくることを、円の上を回る点の動きから確かめます。

\clearpage
%% ---------------- やってみよう・思い出そう ----------------
\begin{yattemiyou}[表計算で $\sin x$ の傾きを調べよう]
$h=0.001$ として、$\dfrac{\sin(x+h)-\sin x}{h}$ を計算します（5-3の「平均の変化率で $h$ を小さくする」と同じ）。$x$ はラジアン（4-2）です。表計算の \texttt{SIN} は、ラジアンで計算します。\bkidx{らじあん}{ラジアン}
\begin{bkcenter}
\bksheet{4}{22}{%
1 & $x$ & $\sin x$ & 傾き（$h=0.001$） & $\cos x$ \\
2 & 0 & 0 & 1.0000 & 1 \\
3 & $\frac{\pi}{6}=0.5236$ & 0.5 & 0.8658 & 0.8660 \\
4 & $\frac{\pi}{4}=0.7854$ & 0.7071 & 0.7068 & 0.7071 \\
5 & $\frac{\pi}{3}=1.0472$ & 0.8660 & 0.4996 & 0.5 \\
6 & $\frac{\pi}{2}=1.5708$ & 1 & $-0.0005$ & 0 \\
7 & $\pi=3.1416$ & 0 & $-1.0000$ & $-1$ \\}
\bktabcaption{傾きの列と、$\cos x$ の列がほとんど同じになる。}
\end{bkcenter}
\end{yattemiyou}

\noindent $\sin x$ の傾きは、$\cos x$ にとても近い値になりました。小さなずれ（0.8658 と 0.8660 など）は、$h$ が0でない分です。$\sin x$ を微分すると $\cos x$ になりそうです。この章では、そのわけを円運動で考えます。

\begin{omoidasou}
\bkrecall{4-2}{弧度法：角 $\theta$（rad）＝弧の長さ÷半径。半径1の円なら、角と弧の長さは同じ数。1周は $2\pi$\,rad（360°）。円の接線は、半径に垂直。}
\bkrecall{4-8}{半径1の円の上で、$x$ 軸から角 $\theta$ の所にある点の座標が $(\cos\theta,\ \sin\theta)$。\\半径 $r$ の円の上を、1秒に $\omega$ ラジアン（角速度 $\omega$）で回る点は\\$(r\cos\omega t,\ r\sin\omega t)$。周期 $T=\dfrac{2\pi}{\omega}$}
\bkrecall{4-11}{速度は、向きと大きさをもつ量（ベクトル）。$x$ 方向の成分と $y$ 方向の成分に分けて考えられる。}
\bkrecall{5-5}{速度は位置を時間で微分したもの、加速度は速度を微分したもの。力＝質量×加速度。}
\bkrecall{5-6}{合成関数の微分：外側を微分して、内側の微分をかける。}
\end{omoidasou}

\clearpage
%% ---------------- 数学の道具 ----------------
\bksection{doguu}{数学の道具}

\bksubsection{1}{円の上を回る点の速さと向き}

\label{bk:5-15:en}半径1の円の上を、点Pが「1秒に1ラジアン」の割合で回ります。0秒のとき、Pは $x$ 軸の上（角0）にいます。$t$ 秒たつと、Pは角 $t$ の所にいるので、Pの座標は $(\cos t,\ \sin t)$ です（4-8）。

\noindent\begin{minipage}[c]{0.46\linewidth}
\centering
\begin{tikzpicture}[x=1mm, y=1mm]
  \draw[bkGray, line width=0.6pt, -{Stealth[length=4pt]}] (-26,0) -- (28,0) node[anchor=north, font=\sffamily\scriptsize] {$x$};
  \draw[bkGray, line width=0.6pt, -{Stealth[length=4pt]}] (0,-26) -- (0,34) node[anchor=east, font=\sffamily\scriptsize] {$y$};
  \draw[bkBlue, line width=1.2pt] (0,0) circle (21);
  \def\ang{35}
  \draw[bkInk, line width=0.8pt] (0,0) -- (\ang:21);
  \fill[bkInk] (\ang:21) circle (1.1);
  \draw[bk arrow, bkOrange, line width=1.5pt] (\ang:21) -- ++(\ang+90:15);
  \draw[bkInk, line width=0.5pt] (0,0) ++(0:6) arc (0:\ang:6);
  \node[font=\sffamily\scriptsize] at ({\ang/2}:8.5) {$t$};
  \node[font=\sffamily\scriptsize, anchor=west] at ($(\ang:21)+(1,-1.5)$) {P};
  \node[font=\sffamily\scriptsize, text=bkOrange!80!black, anchor=south west, fill=white, inner sep=1pt] at ($(\ang:21)+(\ang+90:15)+(1,-1)$) {速度（長さ1）};
  \draw[bkInk, line width=0.4pt] (\ang:18) -- ++(\ang+90:3) -- ++(\ang:3);
\end{tikzpicture}
\end{minipage}\hfill
\begin{minipage}[c]{0.5\linewidth}
\begin{itemize}[leftmargin=1.2\zw, itemsep=3pt, topsep=0pt, label={\color{bkTheme}●}]
\item \textbf{速さ}\quad 弧の長さ＝半径×角（4-2）です。1秒で角は1ふえるので、1秒に進む弧の長さは $1\times1=1$。だから\textbf{速さは1}です。
\item \textbf{向き}\quad Pは円にそって動くので、向きは円の接線の向きです。接線は半径に垂直です（4-2）。だから速度の矢印は、\textbf{半径の向きを左回りに90°回した向き}です（Pは左回りに回っているから）。
\end{itemize}
\end{minipage}\par\vspace{4pt}
この「長さ1で、半径に垂直な矢印」を、$x$ 方向（横）と $y$ 方向（たて）の2つに分けてみましょう（4-11）。

\bksubsection{2}{速度の矢印を、たてと横に分ける}

\label{bk:5-15:bunkai}下の図の左は、半径（長さ1）を、横 $\cos t$ とたて $\sin t$ の直角三角形の斜めの辺と見たものです。速度の矢印は、この半径を90°回したものです。だから、\textbf{この直角三角形をまるごと90°回した三角形}が、速度の矢印の下にできます。

\begin{bkfigure}
\begin{tikzpicture}[x=1mm, y=1mm]
  \def\R{30}
  \draw[bkGray, line width=0.6pt, -{Stealth[length=4pt]}] (-6,0) -- (40,0) node[anchor=west, font=\sffamily\scriptsize] {$x$};
  \draw[bkGray, line width=0.6pt, -{Stealth[length=4pt]}] (0,-4) -- (0,46) node[anchor=east, font=\sffamily\scriptsize] {$y$};
  \draw[bkBlue!50, line width=0.8pt] (-4:\R) arc (-4:48:\R);
  % 半径の三角形（横 cos t は青、たて sin t はオレンジ）
  \fill[bkBlue!8] (0,0) -- (25.98,0) -- (25.98,15) -- cycle;
  \draw[bkBlue, line width=1.6pt] (0,0) -- (25.98,0);
  \draw[bkOrange, line width=1.6pt] (25.98,0) -- (25.98,15);
  \draw[bkInk, line width=0.9pt] (0,0) -- (25.98,15);
  \fill[bkInk] (25.98,15) circle (1);
  \node[font=\sffamily\scriptsize, anchor=north] at (13,-1.5) {$\cos t$};
  \node[font=\sffamily\scriptsize, anchor=east] at (24.4,4.6) {$\sin t$};
  \draw[bkInk, line width=0.5pt] (7,0) arc (0:30:7);
  \node[font=\sffamily\scriptsize] at (15:9.8) {$t$};
  \node[font=\sffamily\scriptsize, anchor=north west] at (26.8,15) {P};
  % 速度の三角形（90°回した。横 sin t はオレンジ、たて cos t は青）
  \fill[bkOrange!8] (25.98,15) -- (10.98,15) -- (10.98,40.98) -- cycle;
  \draw[bkOrange, line width=1.6pt] (25.98,15) -- (10.98,15);
  \draw[bkBlue, line width=1.6pt] (10.98,15) -- (10.98,40.98);
  \draw[bk arrow, bkInk, line width=1.1pt] (25.98,15) -- (10.98,40.98);
  \node[font=\sffamily\scriptsize, anchor=south] at (17.5,15.6) {$\sin t$};
  \node[font=\sffamily\scriptsize, anchor=east] at (9.4,28) {$\cos t$};
  \node[font=\sffamily\scriptsize, anchor=west, fill=white, inner sep=1pt] at (20.5,30) {速度（長さ1）};
  % 右の説明
  \node[font=\sffamily\small, anchor=north west, align=left, text width=66mm] at (52,44) {%
  \textcolor{bkBlue!80!black}{青い辺}の長さは、どちらも $\cos t$。\\
  \textcolor{bkOrange!80!black}{オレンジの辺}の長さは、どちらも $\sin t$。\\[4pt]
  速度の矢印は、\\
  \quad 横に $\sin t$ だけ\textbf{左}（$x$ がへる向き）\\
  \quad たてに $\cos t$ だけ\textbf{上}（$y$ がふえる向き）\\[4pt]
  だから、速度の成分は\\
  \quad $(x\text{ 方向},\ y\text{ 方向})=(-\sin t,\ \cos t)$};
\end{tikzpicture}
\bkfigcaption{半径の三角形（青い面）を90°回すと、速度の三角形（オレンジの面）になる。辺の長さはそのまま、たてと横が入れかわる。}
\end{bkfigure}

\bkvomidashi{数で確かめる}
$t=\dfrac{\pi}{6}$（30°）のとき、$\cos\dfrac{\pi}{6}=0.866$、$\sin\dfrac{\pi}{6}=0.5$ なので、Pは $(0.866,\ 0.5)$ にいます。速度の成分は $(-0.5,\ 0.866)$ のはずです。

ほんとうにそうか、0.001秒だけ時間を進めて調べます。
\begin{bkcenter}
\begin{bktable}{colspec={Q[l,m,wd=34mm] Q[c,m,wd=30mm] Q[c,m,wd=30mm]}, rowsep=1.5pt}
 & $x=\cos t$ & $y=\sin t$ \\
$t=\frac{\pi}{6}$ & 0.866025 & 0.500000 \\
$t=\frac{\pi}{6}+0.001$ & 0.865525 & 0.500866 \\
変化÷0.001 & $-0.500$ & $0.866$ \\
\end{bktable}
\bktabcaption{0.001秒のあいだの変化を0.001でわると、1秒あたりの変化（速度）になる。}
\end{bkcenter}
$x$ は1秒あたり0.5の割合でへり、$y$ は1秒あたり0.866の割合でふえています。図から考えた $(-\sin t,\ \cos t)=(-0.5,\ 0.866)$ とぴったり合います。

\bksubsection{3}{座標を微分すると、速度の成分}

\label{bk:5-15:bibun}5-5で、まっすぐ動く物の速度は「位置を時間で微分したもの」でした。平面の上を動く点でも、横の動きとたての動きを別々に見れば、同じことがいえます。
\begin{itemize}[leftmargin=1.2\zw, itemsep=2pt, topsep=2pt, label={\color{bkTheme}●}]
\item Pの $x$ 座標（横の位置）がふえる速さ＝速度の $x$ 成分
\item Pの $y$ 座標（たての位置、高さ）がふえる速さ＝速度の $y$ 成分
\end{itemize}
Pの $x$ 座標は $\cos t$、$y$ 座標は $\sin t$ でした。そして、速度の成分は $(-\sin t,\ \cos t)$ でした。2つを組にして読むと、
\begin{bkcenter}
\begin{bktable}{colspec={Q[l,m,wd=30mm] Q[c,m,wd=30mm] Q[c,m,wd=50mm]}, rowsep=2pt}
 & 位置（座標） & 位置を微分したもの（速度） \\
$x$ 方向（横） & $\cos t$ & $-\sin t$ \\
$y$ 方向（たて） & $\sin t$ & $\cos t$ \\
\end{bktable}
\end{bkcenter}
つまり、
\[ (\sin t)'=\cos t,\qquad (\cos t)'=-\sin t \]
$\cos t$ の微分に $-$ がつくのは、Pが左回りに上がっていくとき（$0<t<\frac{\pi}{2}$）、横の位置 $x$ がへっていくからです。\bkidx{さいん・こさいんのびぶん}{$\sin$・$\cos$ の微分}

\bkvomidashi{\bkhatenaref の答え}
観覧車のゴンドラの高さが $\sin t$ で表されるとき、高さの変わる速さは $\cos t$ です。$\cos t$ の大きさがいちばん大きい（1 か $-1$）のは、$t=0$ や $t=\pi$、つまりゴンドラが中心と同じ高さ（真横）を通るときです。このとき、速度の矢印はまっすぐ上か下を向いていて、速さ1がそのまま高さの変化になります。いちばん上（$t=\frac{\pi}{2}$）やいちばん下では $\cos t=0$ で、速度の矢印は真横を向いているので、高さは一瞬変わりません。

\noindent\begin{minipage}[c]{0.57\linewidth}
\textbf{いちばん下からはかった高さ}\par
観覧車のゴンドラの高さを、いちばん下からはかってみます。半径\qty{50}{m}の観覧車の中心は、いちばん下から\qty{50}{m}の高さです。いちばん下から角 $\theta$ だけ回ったゴンドラは、中心より $50\cos\theta$ だけ下にいます（右の図の直角三角形）。だから高さは
\[ y=50-50\cos\theta \]
$\theta=0$（いちばん下）で $y=0$、$\theta=\pi$（いちばん上）で $y=50+50=100$。合っています。
\end{minipage}\hfill
\begin{minipage}[c]{0.4\linewidth}
\centering
\begin{tikzpicture}[x=1mm, y=1mm]
  \draw[bkGray, line width=0.8pt] (-26,0) -- (24,0);
  \draw[bkInk!50, line width=0.8pt] (0,20) circle (20);
  \draw[bkGray, dashed, line width=0.5pt] (0,20) -- (0,0);
  \def\th{50}
  \draw[bkInk, line width=0.8pt] (0,20) -- ({20*sin(\th)},{20-20*cos(\th)});
  \draw[bkOrange, line width=1.4pt] (0,20) -- (0,{20-20*cos(\th)});
  \draw[bkGray, line width=0.5pt] (0,{20-20*cos(\th)}) -- ({20*sin(\th)},{20-20*cos(\th)});
  \fill[bkRoute2] ({20*sin(\th)},{20-20*cos(\th)}) ++(-1.8,-2.6) rectangle ++(3.6,3);
  \draw[bkInk, line width=0.5pt] (0,14) arc (-90:{-90+\th}:6);
  \node[font=\sffamily\scriptsize] at ($(0,20)+({-90+\th/2}:8.8)$) {$\theta$};
  \node[font=\sffamily\scriptsize, text=bkOrange!80!black, anchor=east] at (-1.2,13.5) {$50\cos\theta$};
  \draw[bkGray, line width=0.5pt, <->] (-22,0) -- (-22,20);
  \node[font=\sffamily\scriptsize, anchor=east] at (-22.8,10) {50};
  \fill[bkInk] (0,20) circle (0.8);
\end{tikzpicture}
\end{minipage}
例題6では、この式を使って、上がる速さを計算します。


\bksubsection{4}{微分すると、波が $\frac{1}{4}$ 周期ずれる}

\label{bk:5-15:zure}$\sin t$ と $\cos t$ のグラフを重ねてみます。

\begin{bkfigure}
\begin{tikzpicture}[tl/.style={font=\sffamily\footnotesize, text=bkInk, anchor=base, yshift=-6mm}]
\begin{axis}[bkplot, width=120mm, height=46mm, xmin=0, xmax=7, ymin=-1.15, ymax=1.35, xtick={1.5708,3.1416,4.7124,6.2832}, xticklabels={}, ytick={-1,1}, xlabel={$t$}, ylabel={$y$}]
  \addplot[bkcurve, domain=0:6.6, samples=120] {sin(deg(x))};
  \addplot[bkcurve2, domain=0:6.6, samples=120] {cos(deg(x))};
  % 目もりの数字は、波と重ならないように −1 の線より下に置く。分数も π・2π と同じ字の大きさにし、ベースラインをそろえる
  \node[tl] at (axis cs:1.5708,-1.15) {$\dfrac{\pi}{2}$};
  \node[tl] at (axis cs:3.1416,-1.15) {$\pi$};
  \node[tl] at (axis cs:4.7124,-1.15) {$\dfrac{3\pi}{2}$};
  \node[tl] at (axis cs:6.2832,-1.15) {$2\pi$};
  \draw[bk arrow, bkGray, line width=0.8pt] (axis cs:0,1.15) -- (axis cs:1.5708,1.15);
  \node[font=\sffamily\scriptsize, anchor=south] at (axis cs:0.785,1.15) {$\frac{1}{4}$ 周期};
  \node[font=\sffamily\scriptsize, text=bkBlue, anchor=west] at (axis cs:6.75,0.31) {$\sin t$};
  \node[font=\sffamily\scriptsize, text=bkOrange!80!black, anchor=west] at (axis cs:6.75,0.95) {$\cos t$};
\end{axis}
\end{tikzpicture}
\bkfigcaption{$\sin t$（青）の傾きのグラフが $\cos t$（オレンジ）。$\cos t$ の山は $t=0$、$\sin t$ の山は $t=\frac{\pi}{2}$。$\cos t$ のほうが $\frac{\pi}{2}$ だけ早く山が来る。}
\end{bkfigure}
\begin{itemize}[leftmargin=1.2\zw, itemsep=2pt, topsep=2pt, label={\color{bkTheme}●}]
\item $\sin t$ が0を上向きに通る所（$t=0$）では、青い線がいちばん急に上がっていて、傾きは最大の1。そこで $\cos t$ はちょうど山（1）です。
\item $\sin t$ が山（$t=\frac{\pi}{2}$）では、青い線は一瞬平らで、傾きは0。そこで $\cos t=0$ です。
\item $\sin t$ が0を下向きに通る所（$t=\pi$）では、傾きは最も急な下り（$-1$）。そこで $\cos t$ は谷（$-1$）です。
\end{itemize}
$\cos t$ は、$\sin t$ と同じ形の波を、$\frac{\pi}{2}$ だけ左にずらしたものです。式で書くと $\cos t=\sin\bigl(t+\frac{\pi}{2}\bigr)$。数で確かめましょう。
\begin{bkcenter}
\begin{bktable}{colspec={Q[l,m,wd=30mm] *{4}{Q[c,m,wd=18mm]}}, rowsep=1.5pt}
$t$ & 0 & $\dfrac{\pi}{6}$ & $\dfrac{\pi}{3}$ & $\dfrac{\pi}{2}$ \\
$\cos t$ & 1 & 0.866 & 0.5 & 0 \\
$t+\dfrac{\pi}{2}$ & $\dfrac{\pi}{2}$ & $\dfrac{2\pi}{3}$ & $\dfrac{5\pi}{6}$ & $\pi$ \\
$\sin\Bigl(t+\dfrac{\pi}{2}\Bigr)$ & 1 & 0.866 & 0.5 & 0 \\
\end{bktable}
\end{bkcenter}
1周期は $2\pi$ なので、$\frac{\pi}{2}$ は $2\pi$ の $\frac{1}{4}$ です。\textbf{$\sin t$ を微分すると、同じ形の波が $\frac{1}{4}$ 周期だけ先に進みます}。

微分するたびに $\frac{1}{4}$ 周期ずつ進むので、4回微分すると1周期進んで、もとの $\sin t$ にもどります。

\begin{bkfigure}
\begin{tikzpicture}[x=1mm, y=1mm, nd/.style={draw=bkTheme, line width=0.9pt, rounded corners=3pt, fill=bkTheme!8, minimum width=20mm, minimum height=8mm, font=\small}]
  \node[nd] (a) at (0,13) {$\sin t$};
  \node[nd] (b) at (34,0) {$\cos t$};
  \node[nd] (c) at (0,-13) {$-\sin t$};
  \node[nd] (d) at (-34,0) {$-\cos t$};
  \draw[bk arrow, bkOrange, line width=1.1pt] (a.east) to[bend left=20] node[pos=0.5, anchor=south west, font=\sffamily\scriptsize] {微分} (b.north);
  \draw[bk arrow, bkOrange, line width=1.1pt] (b.south) to[bend left=20] node[pos=0.5, anchor=north west, font=\sffamily\scriptsize] {微分} (c.east);
  \draw[bk arrow, bkOrange, line width=1.1pt] (c.west) to[bend left=20] node[pos=0.5, anchor=north east, font=\sffamily\scriptsize] {微分} (d.south);
  \draw[bk arrow, bkOrange, line width=1.1pt] (d.north) to[bend left=20] node[pos=0.5, anchor=south east, font=\sffamily\scriptsize] {微分} (a.west);
  \node[font=\sffamily\scriptsize, text=bkGray, align=center] at (0,0) {右回り：微分\\左回り：積分};
\end{tikzpicture}
\bkfigcaption{微分すると右回りに1つ進む。4回でもとにもどる。積分（$+C$ は省略）は左回り。}
\end{bkfigure}
図の3つめと4つめは、1つめと2つめに $-$ をつけたものです。定数倍（ここでは $-1$ 倍）はそのまま前に出る（5-4）ので、
\begin{align*}
(-\sin t)'&=-(\sin t)'=-\cos t\\
(-\cos t)'&=-(\cos t)'=-(-\sin t)=\sin t
\end{align*}

\bksubsection{5}{角速度 $\omega$ をふくむとき}

\label{bk:5-15:omega}\bkidx{かくそくど}{角速度}1秒に2ラジアン回る点を考えます。$t$ 秒後の角は $2t$ なので、高さは $\sin2t$ です。2倍の速さで回るので、高さも2倍の速さで変わるはずです。

合成関数の微分（5-6）で確かめます。内側を $u=2t$、外側を $\sin u$ とすると、
\[ (\sin2t)'=\underbrace{\cos u}_{\text{外側の微分}}\times\underbrace{(2t)'}_{\text{内側の微分}}=\cos2t\times2=2\cos2t \]
前に出た2は、「1秒に2ラジアン回る」の2です。表計算でも確かめましょう。$t=0.3$ では、
\begin{align*}
\text{傾き（$h=0.001$）}&\quad \frac{\sin(2\times0.301)-\sin(2\times0.3)}{0.001}=1.6495\\
\text{微分の式}&\quad 2\cos(2\times0.3)=2\cos0.6=1.6507
\end{align*}
で、ほとんど同じです。

1秒に $\omega$ ラジアン回るときは、角が $\omega t$ です。同じように内側を $u=\omega t$（微分すると $\omega$）として、
\[ (\sin\omega t)'=\omega\cos\omega t,\qquad (\cos\omega t)'=-\omega\sin\omega t \]
$\omega$（オメガ）は\textbf{角速度}といい、単位は rad/s です。

\bkvomidashi{半径が $r$ のとき}
半径 $r$ の円では、座標はどちらも半径1のときの $r$ 倍です。$r$ は定数なので、微分しても前にそのまま残ります（5-4）。
\[ \text{座標}\ (r\cos\omega t,\ r\sin\omega t)\quad\xrightarrow{\ \text{微分}\ }\quad\text{速度}\ (-r\omega\sin\omega t,\ r\omega\cos\omega t) \]
速さ（速度の矢印の長さ）は $r\omega$ です。1秒に回る角が $\omega$ なので、1秒に進む弧の長さは「半径×角」で $r\omega$（4-2）。微分した結果と合っています。速く回る（$\omega$ が大きい）ほど、大きい円（$r$ が大きい）ほど、速く動きます。

\bksubsection{6}{2回微分すると、もとの $-\omega^2$ 倍}

\label{bk:5-15:nikai}$\sin\omega t$ を2回微分してみます。1回めは⑤のとおり $\omega\cos\omega t$。2回めは、$\omega$ を前に置いたまま $\cos\omega t$ を微分します。
\[ (\sin\omega t)''=(\omega\cos\omega t)'=\omega\times(-\omega\sin\omega t)=-\omega^2\sin\omega t \]
$\omega$ が2回かかって $\omega^2$、$\cos$ の微分で $-$ がついて、$-\omega^2$ になりました。$\cos\omega t$ でも同じで、\textbf{2回微分すると、もとの関数の $-\omega^2$ 倍}になります。ゆれの位置を $x=A\sin\omega t$（$A$ はゆれの大きさ）とすると、$\ddot{x}$（$x$ の2回微分、5-5）は
\[ \ddot{x}=-\omega^2x \]
\textbf{加速度が、位置に比例して、いつも位置と反対向き}です。$x$ が正（中心より右）なら加速度は負（左向き）、$x$ が負なら加速度は正（右向き）。いつも中心のほうへ引きもどされます。

\bkvomidashi{ばねのおもり}
ばねは、のびに比例する力（2-15の $F=kx$）で、のびと反対向きにおもりを引きもどします。向きまでこめて書くと $F=-kx$。力＝質量×加速度（5-5）なので、
\[ m\ddot{x}=-kx\qquad\text{両辺を }m\text{ でわると}\qquad \ddot{x}=-\frac{k}{m}\,x \]
これは $\omega^2=\dfrac{k}{m}$ とした $\ddot{x}=-\omega^2x$ と同じ形です。だから、ばねのおもりは $\sin$ の形でゆれます。このゆれを\textbf{単振動}といいます。たとえば質量\qty{0.5}{kg}のおもりと、ばね定数\qty{8}{N/m}のばねなら、$\omega^2=8\div0.5=16$ で $\omega=4$、周期は $\dfrac{2\pi}{4}\approx\qty{1.57}{s}$ です。くわしくは6-6で調べます。\bkidx{たんしんどう}{単振動}

\begin{bkfigure}
\begin{tikzpicture}[tl/.style={font=\sffamily\footnotesize, text=bkInk, anchor=base, yshift=-6mm}]
\begin{axis}[bkplot, width=120mm, height=46mm, xmin=0, xmax=6.6, ymin=-1.15, ymax=1.5, xtick={1.5708,3.1416,4.7124,6.2832}, xticklabels={}, ytick={-1,1}, xlabel={$t$}]
  \addplot[bkcurve, domain=0:6.4, samples=120] {sin(deg(x))};
  \addplot[bkcurve2, domain=0:6.4, samples=120] {cos(deg(x))};
  \addplot[bkcurve3, domain=0:6.4, samples=120] {-sin(deg(x))};
  \node[font=\sffamily\scriptsize, text=bkBlue, anchor=south west, inner sep=1pt] at (axis cs:1.7,1.12) {位置 $x$};
  \node[font=\sffamily\scriptsize, text=bkOrange!80!black, anchor=south west, inner sep=1pt] at (axis cs:0.12,1.12) {速度 $v$};
  \node[font=\sffamily\scriptsize, text=bkGreen!70!black, anchor=south west, inner sep=1pt] at (axis cs:4.84,1.12) {加速度 $a$};
  \draw[bkGray, dashed, line width=0.5pt] (axis cs:1.5708,-1.1) -- (axis cs:1.5708,1.0);
  \draw[bkGray, dashed, line width=0.5pt] (axis cs:3.1416,-1.1) -- (axis cs:3.1416,1.0);
  % 目もりの数字と「はし」「中心」は、波と重ならないように −1 の線より下に置く
  \node[tl] at (axis cs:1.5708,-1.15) {$\dfrac{\pi}{2}$};
  \node[tl] at (axis cs:3.1416,-1.15) {$\pi$};
  \node[tl] at (axis cs:4.7124,-1.15) {$\dfrac{3\pi}{2}$};
  \node[tl] at (axis cs:6.2832,-1.15) {$2\pi$};
  \node[tl, font=\sffamily\scriptsize, yshift=-6.5mm] at (axis cs:1.5708,-1.15) {はし};
  \node[tl, font=\sffamily\scriptsize, yshift=-6.5mm] at (axis cs:3.1416,-1.15) {中心};
\end{axis}
\end{tikzpicture}
\bkfigcaption{$x=\sin t$ のとき、$v=\cos t$、$a=-\sin t$。はし（$x$ が最大）では速度0で加速度が最大、中心では速度が最大で加速度0。}
\end{bkfigure}
図の緑の線（加速度）は、青い線（位置）を上下にひっくり返した形です。これが「$-1$ 倍」（ここでは $\omega=1$ なので $-\omega^2=-1$）の意味です。

\bksubsection{7}{$\sin$・$\cos$ の積分}

\label{bk:5-15:sekibun}積分は微分の逆でした（5-11）。微分の式を、右から左へ逆に読みます。
\begin{bkcenter}
\begin{bktable}{colspec={Q[c,m,wd=40mm] Q[c,m,wd=58mm]}, rowsep=2pt}
微分の式 & 逆に読んだ積分の式 \\
$(-\cos t)'=\sin t$ & $\displaystyle\int\sin t\,\dd t=-\cos t+C$ \\
$(\sin t)'=\cos t$ & $\displaystyle\int\cos t\,\dd t=\sin t+C$ \\
\end{bktable}
\end{bkcenter}
$\sin t$ の積分に $-$ がつくのは、$\cos t$ を微分すると $-\sin t$ になるからです。$-\cos t$ を微分すれば、$-$ が2つで $+\sin t$ になります。

\bkvomidashi{角速度 $\omega$ があるとき}
$\cos\omega t$ を微分すると $-\omega\sin\omega t$ で、ほしい $\sin\omega t$ の $-\omega$ 倍になってしまいます。そこで、はじめに $-\dfrac{1}{\omega}$ をかけておくと、微分したとき $\omega$ と打ち消しあいます。
\[ \Bigl(-\frac{1}{\omega}\cos\omega t\Bigr)'=-\frac{1}{\omega}\times(-\omega\sin\omega t)=\sin\omega t \]
だから、
\[ \int\sin\omega t\,\dd t=-\frac{1}{\omega}\cos\omega t+C,\qquad \int\cos\omega t\,\dd t=\frac{1}{\omega}\sin\omega t+C \]
\textbf{微分では $\omega$ をかけ、積分では $\omega$ でわる}。微分と積分は逆のはたらきなので、$\omega$ のあつかいも逆になります。

\bkvomidashi{例}
$\sin t$ の山1つ分（0から $\pi$ まで）の面積を求めます。$-\cos t$ の、上端 $\pi$ での値から下端0での値をひきます（5-11）。$\cos\pi=-1$、$\cos0=1$ なので、
\begin{align*}
\int_0^\pi\sin t\,\dd t&=\Bigl[\,-\cos t\,\Bigr]_0^\pi=(-\cos\pi)-(-\cos0)\\
&=\{-(-1)\}-(-1)=1+1=2
\end{align*}
5-10のように、細い長方形をたしても確かめられます。0から $\pi$ を $n$ 等分し、まん中の高さで長方形をつくってたすと、$n=6$ で2.023、$n=12$ で2.006、$n=100$ で2.0001。2に近づきます。

\begin{bkfigure}
\begin{tikzpicture}
\begin{axis}[bkplot, name=L, width=64mm, height=36mm, xmin=0, xmax=3.6, ymin=0, ymax=1.3, xtick={1.5708,3.1416}, xticklabels={$\dfrac{\pi}{2}$,$\pi$}, ytick={1}, xlabel={$t$}, title={\small $\sin t$ の山1つ分}, title style={yshift=-1mm}]
  \addplot[fill=bkBlue!15, draw=none, domain=0:3.1416, samples=80] {sin(deg(x))} \closedcycle;
  \addplot[bkcurve, domain=0:3.1416, samples=80] {sin(deg(x))};
  \node[font=\sffamily\scriptsize, text=bkBlue!80!black] at (axis cs:1.5708,0.4) {面積2};
\end{axis}
\begin{axis}[bkplot, at={(L.east)}, anchor=west, xshift=10mm, width=64mm, height=36mm, xmin=0, xmax=3.6, ymin=0, ymax=1.3, xtick={1.5708,3.1416}, xticklabels={$\dfrac{\pi}{2}$,$\pi$}, ytick={1}, xlabel={$t$}, title={\small $\sin2t$ の山1つ分}, title style={yshift=-1mm}]
  \addplot[fill=bkOrange!25, draw=none, domain=0:1.5708, samples=60] {sin(deg(2*x))} \closedcycle;
  \addplot[bkcurve2, domain=0:1.5708, samples=60] {sin(deg(2*x))};
  \node[font=\sffamily\scriptsize, text=bkOrange!80!black] at (axis cs:0.785,0.4) {面積1};
\end{axis}
\end{tikzpicture}
\bkfigcaption{$\sin2t$（オレンジ）の山は、$\sin t$（青）の山と高さは同じで、はばが半分。だから面積も半分の1。積分で $\frac{1}{\omega}=\frac{1}{2}$ がかかるのは、このため。}
\end{bkfigure}
$\sin2t$ の山1つ分は0から $\frac{\pi}{2}$ までです（$2t$ が0から $\pi$ まで動くから）。計算すると、
\[ \int_0^{\frac{\pi}{2}}\sin2t\,\dd t=\Bigl[\,-\frac{1}{2}\cos2t\,\Bigr]_0^{\frac{\pi}{2}}=\Bigl(-\frac{1}{2}\cos\pi\Bigr)-\Bigl(-\frac{1}{2}\cos0\Bigr)=\frac{1}{2}+\frac{1}{2}=1 \]
図の「はばが半分なので面積も半分」と合います。

\bksubsection{8}{なぜラジアンを使うのか}

\label{bk:5-15:rad}角を度（°）ではかると、どうなるでしょう。30°から1°ふやしたときの $\sin$ のふえ方は
\[ \sin31^\circ-\sin30^\circ=0.51504-0.50000=0.01504 \]
で、$\cos30^\circ=0.866$ とは合いません。1°は $\dfrac{\pi}{180}$ ラジアン（$\pi$ ラジアンが180°だから）なので、1°あたりのふえ方は $0.866\times\dfrac{\pi}{180}=0.0151$ になります。度で微分すると、いつも $\dfrac{\pi}{180}$ という余分な数がかかるのです。

ラジアンで角をはかると、半径1の円で「角＝弧の長さ」になり、①で見たように速さがちょうど1になるので、$(\sin t)'=\cos t$ ときれいな形になります。微分や積分では、角はいつもラジアンで表します。電卓や表計算で $\sin$ を計算するときも、角がラジアンになっているか確かめましょう。

\bksubsection{9}{$\sin^2$ の平均：交流の電力}

\label{bk:5-15:jikko}\bkvomidashi{交流の電力は、行ったり来たりする}
家のコンセントの電圧は、$V=V_0\sin\omega t$ のように $+$ と $-$ を行ったり来たりする交流です（$V_0$ は山の高さ）。抵抗 $R$ をつなぐと、電流は $I=\dfrac{V}{R}$（2-15）。電力は電圧×電流なので、
\[ P=VI=V\times\frac{V}{R}=\frac{V^2}{R}=\frac{V_0^2}{R}\sin^2\omega t \]
$\sin^2\omega t$ は $(\sin\omega t)^2$ のことです。2乗なので、電圧が $-$ のときも電力は $+$ で、0から $\dfrac{V_0^2}{R}$ のあいだを行ったり来たりします。では、平均するといくらでしょう。

\bkvomidashi{数で見当をつける}
$\sin^2t$ の値を、1周期（0から $2\pi$）のあいだ、$\frac{\pi}{6}$（30°）ごとに12こ調べます。
\begin{bkcenter}
\begin{bktable}{colspec={Q[l,m,wd=16mm] *{6}{Q[c,m,wd=13mm]}}, rowsep=1.5pt}
角 & 0° & 30° & 60° & 90° & 120° & 150° \\
$\sin t$ & 0 & 0.5 & 0.87 & 1 & 0.87 & 0.5 \\
$\sin^2t$ & 0 & 0.25 & 0.75 & 1 & 0.75 & 0.25 \\
\end{bktable}\par\vspace{4pt}
\begin{bktable}{colspec={Q[l,m,wd=16mm] *{6}{Q[c,m,wd=13mm]}}, rowsep=1.5pt}
角 & 180° & 210° & 240° & 270° & 300° & 330° \\
$\sin t$ & 0 & $-0.5$ & $-0.87$ & $-1$ & $-0.87$ & $-0.5$ \\
$\sin^2t$ & 0 & 0.25 & 0.75 & 1 & 0.75 & 0.25 \\
\end{bktable}
\end{bkcenter}
後半の6こは、$\sin t$ が $-$ になっても、2乗すると前半と同じ値になります。$\sin^2t$ の12この合計は $(0+0.25+0.75+1+0.75+0.25)\times2=6$ なので、平均は $6\div12=0.5$。どうやら平均は $\frac{1}{2}$ になりそうです。0.25と0.75、0と1のように、平均より小さい値と大きい値が、ちょうど組になっているからです。

\noindent\begin{minipage}[c]{0.5\linewidth}
\begin{tikzpicture}
\begin{axis}[bkplot, width=66mm, height=36mm, xmin=0, xmax=6.6, ymin=0, ymax=1.15, xtick={3.1416,6.2832}, xticklabels={$\pi$,$2\pi$}, ytick={0.5,1}, yticklabels={$\dfrac{1}{2}$,1}, xlabel={$t$}]
  \addplot[bkcurve, domain=0:6.2832, samples=120, fill=bkBlue!12] {sin(deg(x))^2} \closedcycle;
  \addplot[bkcurve2, domain=0:6.2832] {0.5};
\end{axis}
\end{tikzpicture}
\end{minipage}\hfill
\begin{minipage}[c]{0.46\linewidth}
グラフで見ると、$\sin^2t$（青）の山の、オレンジの線より上の部分を切りとって、谷にうめると、ちょうど高さ $\frac{1}{2}$ の平らな形になります。平均が $\frac{1}{2}$ というのは、このことです。
\end{minipage}\par\vspace{4pt}
これを積分で確かめます。そのために、$\sin^2$ を、積分のできる形に変えます。

\bkvomidashi{2倍角の式で、$\sin^2$ を $\cos2\theta$ に変える}
4-9の加法定理 $\cos(a+b)=\cos a\cos b-\sin a\sin b$ で、$a$ も $b$ も $\theta$ にすると、
\[ \cos2\theta=\cos\theta\cos\theta-\sin\theta\sin\theta=\cos^2\theta-\sin^2\theta \]
$\sin^2\theta+\cos^2\theta=1$（4-8。半径1の円で三平方の定理）から $\cos^2\theta=1-\sin^2\theta$ なので、これを入れると、
\[ \cos2\theta=(1-\sin^2\theta)-\sin^2\theta=1-2\sin^2\theta \]
これを \textbf{2倍角の式}といいます。$\theta=\frac{\pi}{6}$（30°）で確かめると、左は $\cos60^\circ=0.5$、右は $1-2\times0.5^2=1-0.5=0.5$。合っています。\bkidx{にばいかく}{2倍角}

この式を、$\sin^2\theta$ について解きます（2-14）。$2\sin^2\theta$ を左へ、$\cos2\theta$ を右へ移すと $2\sin^2\theta=1-\cos2\theta$。両辺を2でわって、
\[ \sin^2\theta=\frac{1-\cos2\theta}{2},\qquad\text{同じようにして}\qquad \cos^2\theta=\frac{1+\cos2\theta}{2} \]
（$\cos^2\theta$ のほうは、$\cos2\theta=\cos^2\theta-\sin^2\theta$ の $\sin^2\theta$ に $1-\cos^2\theta$ を入れて作ります。$\cos2\theta=2\cos^2\theta-1$ となるので、$\cos^2\theta$ について解きます。）

\bkvomidashi{1周期で積分する}
$\sin^2t=\dfrac{1-\cos2t}{2}$ を、2つに分けて1つずつ積分します。
\[ \frac{1-\cos2t}{2}=\frac{1}{2}-\frac{1}{2}\cos2t \]
\begin{itemize}[leftmargin=1.2\zw, itemsep=1pt, topsep=2pt, label={\color{bkTheme}●}]
\item $\dfrac{1}{2}$（定数）の積分は $\dfrac{1}{2}t=\dfrac{t}{2}$（5-11）。
\item $\cos2t$ の積分は、⑦で $\omega=2$ として $\dfrac{1}{2}\sin2t$。だから $\dfrac{1}{2}\cos2t$ の積分は、その $\dfrac{1}{2}$ 倍の $\dfrac{1}{4}\sin2t=\dfrac{\sin2t}{4}$。
\end{itemize}
\begin{align*}
\int_0^{2\pi}\sin^2t\,\dd t&=\int_0^{2\pi}\Bigl(\frac{1}{2}-\frac{1}{2}\cos2t\Bigr)\dd t=\Bigl[\,\frac{t}{2}-\frac{\sin2t}{4}\,\Bigr]_0^{2\pi}\\
&=\Bigl(\frac{2\pi}{2}-\frac{\sin4\pi}{4}\Bigr)-\Bigl(0-\frac{\sin0}{4}\Bigr)=(\pi-0)-(0-0)=\pi
\end{align*}
平均は、積分÷区間の長さ（5-12）なので、$\pi\div2\pi=\dfrac{1}{2}$。表で見当をつけた値と合いました。\textbf{$\sin^2$ の1周期の平均は $\frac{1}{2}$}です。$\cos^2$ も同じで、平均は $\frac{1}{2}$ です。

\bkvomidashi{実効値}
だから、交流の平均の電力は $\dfrac{V_0^2}{R}\times\dfrac{1}{2}=\dfrac{V_0^2}{2R}$ です。では、同じ抵抗に直流の電圧 $V_{\mathrm{e}}$ をかけて、同じ平均の電力 $\dfrac{V_{\mathrm{e}}^2}{R}$ を出すには、$V_{\mathrm{e}}$ をいくらにすればよいでしょう。2つの電力を等しいとおくと、
\[ \frac{V_{\mathrm{e}}^2}{R}=\frac{V_0^2}{2R}\qquad\text{両辺に }R\text{ をかけて}\qquad V_{\mathrm{e}}^2=\frac{V_0^2}{2}\qquad\text{つまり}\qquad V_{\mathrm{e}}=\frac{V_0}{\sqrt2} \]
この $V_{\mathrm{e}}$ を、交流の\textbf{実効値}といいます。「はたらき（電力）で見ると、直流の何Vと同じか」を表す値です。\bkidx{じっこうち}{実効値}

家のコンセントの「100\,V」は実効値です。電圧の山の高さは、
\[ V_0=\sqrt2\,V_{\mathrm{e}}=1.414\times100\approx\qty{141}{V} \]
もあります。電気製品は、この山の高さにもたえられるように作られています。

\bksubsection{10}{表計算で、ばねを動かしてみよう}

\label{bk:5-15:sim}5-13では、速度から少しずつ位置を進めて、動きを計算しました。ばねのおもりでも同じことができます。$\omega=1$（$\ddot{x}=-x$）、はじめ $x=1$、$v=0$（はしで手をはなす）とし、$\Delta t=0.1$ 秒ごとに、次の順で計算します。
\begin{enumerate}[label={\bkcircled[bkTheme]{\arabic*}}, leftmargin=2.4\zw, itemsep=1pt, topsep=2pt]
\item 加速度を求める：$a=-x$（⑥の $\ddot{x}=-\omega^2x$ で $\omega=1$）
\item 新しい速度：$v+a\times\Delta t$（加速度は「1秒あたりの速度の変化」だから）
\item 新しい位置：$x+(\text{新しい速度})\times\Delta t$（速度は「1秒あたりの位置の変化」だから）
\end{enumerate}
たとえば1行めは、$a=-1$、新しい速度 $0+(-1)\times0.1=-0.1$、新しい位置 $1+(-0.1)\times0.1=0.99$ です。

\noindent\begin{minipage}[c]{0.44\linewidth}
\centering
\begin{tblr}{colspec={Q[c,m]Q[c,m]Q[c,m]Q[c,m]}, row{1}={bg=bkTheme!12, font=\sffamily\footnotesize}, cells={font=\footnotesize}, hline{1,2,Z}={0.6pt,bkGray}, hline{8}={0.4pt,bkGray,dashed}, colsep=3pt, rowsep=0.5pt}
$t$ & $x$ & $v$ & $\cos t$ \\
0 & 1 & 0 & 1 \\
0.1 & 0.990 & $-0.100$ & 0.995 \\
0.2 & 0.970 & $-0.199$ & 0.980 \\
0.3 & 0.940 & $-0.296$ & 0.955 \\
0.4 & 0.901 & $-0.390$ & 0.921 \\
0.5 & 0.853 & $-0.480$ & 0.878 \\
3.1 & $-1.001$ & $-0.040$ & $-0.999$ \\
6.3 & 0.999 & $-0.019$ & 1.000 \\
\end{tblr}
\end{minipage}\hfill
\begin{minipage}[c]{0.54\linewidth}
\begin{tikzpicture}
\begin{axis}[bkplot, width=70mm, height=40mm, xmin=0, xmax=6.6, ymin=-1.2, ymax=1.2, xtick={3.1416,6.2832}, xticklabels={}, ytick={-1,1}, xlabel={$t$}]
  \addplot[bkcurve, domain=0:6.4, samples=100] {cos(deg(x))};
  % 目もりの数字は、格子の線が通らないように −1 の線より下に置く
  \node[font=\sffamily\footnotesize, text=bkInk, anchor=base, yshift=-4mm] at (axis cs:3.1416,-1.2) {$\pi$};
  \node[font=\sffamily\footnotesize, text=bkInk, anchor=base, yshift=-4mm] at (axis cs:6.2832,-1.2) {$2\pi$};
  \addplot[only marks, mark=*, mark size=1.3pt, bkOrange] coordinates {(0.0,1.000) (0.3,0.940) (0.6,0.797) (0.9,0.582) (1.2,0.315) (1.5,0.020) (1.8,-0.277) (2.1,-0.549) (2.4,-0.772) (2.7,-0.926) (3.0,-0.997) (3.3,-0.979) (3.6,-0.874) (3.9,-0.690) (4.2,-0.445) (4.5,-0.160) (4.8,0.139) (5.1,0.426) (5.4,0.675) (5.7,0.863) (6.0,0.975) (6.3,0.999)};
\end{axis}
\end{tikzpicture}
\end{minipage}\par\vspace{2pt}
{\small オレンジの点が表計算の $x$、青い線が $\cos t$。}\par
「加速度は位置の $-1$ 倍」というきまりだけから計算したのに、点は $\cos t$ の波にぴったりのり、約6.3秒（$2\pi\approx6.28$）で1周してもどってきました。$\ddot{x}=-\omega^2x$ というきまりが、$\sin$・$\cos$ のゆれを生み出していることが分かります。きまり（微分方程式）から動きを求めることは、第6巻の中心です（6-6）。

\begin{motto}[4]
\textbf{加法定理から導く}\par
4-9の加法定理 $\sin(x+h)=\sin x\cos h+\cos x\sin h$ を使うと、
\[ \frac{\sin(x+h)-\sin x}{h}=\sin x\cdot\frac{\cos h-1}{h}+\cos x\cdot\frac{\sin h}{h} \]
$h=0.01$ のとき、表計算で計算すると、
\[ \frac{\sin h}{h}=0.99998,\qquad \frac{\cos h-1}{h}=-0.005 \]
$h$ をもっと0に近づけると、
\[ \frac{\sin h}{h}\to1,\qquad \frac{\cos h-1}{h}\to0 \]
となるので、$(\sin x)'=\sin x\times0+\cos x\times1=\cos x$ です。半径1の円で、小さい角 $h$ の弧の長さ（$h$）と高さ（$\sin h$）がほぼ等しいことが、$\dfrac{\sin h}{h}\to1$ の意味です。
\end{motto}

\begin{senmonsho}[大学の教科書ふうの書き方]
\begin{bkhonbun}
単振動の変位を $x(t)=A\cos(\omega t+\phi)$ とおくと、速度と加速度は
\[ \dot{x}=-A\omega\sin(\omega t+\phi),\qquad \ddot{x}=-A\omega^2\cos(\omega t+\phi)=-\omega^2x \tag{2.3} \]
となる。$A$ を振幅、$\omega$ を角振動数、$\phi$ を初期位相という。\bkidx{しんぷく}{振幅}\bkidx{かくしんどうすう}{角振動数}\bkidx{しょきいそう}{初期位相}
\end{bkhonbun}
\bkyaku{ゆれの位置を $A\cos(\omega t+\phi)$ とする。速度は、これを微分したもの。加速度は、もう一度微分したもので、位置の $-\omega^2$ 倍になる。$\phi$（ファイ）は「$t=0$ で波がどこから始まるか」を決める角。}
\end{senmonsho}

\begin{column}{サインは「弓のつる」から}
$\sin$ の名前は、長い旅をしてきました。1500年ほど前のインドの天文学では、円の弧を弓に、弦（弧の両はしを結ぶ線分）を弓のつるに見立て、その半分の長さを「つるの半分」とよびました。これが半径1の円での $\sin$ です。アラビア語に訳されたときに音だけが写され、のちにヨーロッパでラテン語に訳すとき、にた形の「ふところ・入り江」という言葉と読みまちがえられ、ラテン語の sinus（入り江）になったといわれます。
\end{column}

\bkkigocards{\bkkigocard{$\phi$}{ファイ}{初期位相。$t=0$ での角}{phase}{5-15}}
{\bkkigocard{$V_{\mathrm{e}}$}{ブイ・イー}{実効値。平均の電力が同じになる直流の電圧。$V_0/\sqrt2$}{effective value}{5-15}}

\begin{taisetsu}
\def\tsl#1{\makebox[3.2\zw][l]{\sffamily\bfseries\small\color{bkTheme}#1}}
\tsl{微分}$(\sin\omega t)'=\omega\cos\omega t$\qquad $(\cos\omega t)'=-\omega\sin\omega t$\\[3pt]
\tsl{}2回微分すると、もとの $-\omega^2$ 倍（$\ddot{x}=-\omega^2x$）\\[3pt]
\tsl{積分}$\displaystyle\int\sin\omega t\,\dd t=-\frac{1}{\omega}\cos\omega t+C$\qquad $\displaystyle\int\cos\omega t\,\dd t=\frac{1}{\omega}\sin\omega t+C$\\[3pt]
\tsl{平均}$\sin^2$ の1周期の平均は $\dfrac{1}{2}$\qquad 実効値 $V_{\mathrm{e}}=\dfrac{V_0}{\sqrt2}$\\[3pt]
\tsl{}角は、いつもラジアンではかる。
\end{taisetsu}

\begin{kotoba}
\small
\bkword{角速度}{1秒に回る角 $\omega$（rad/s）。周期は $\frac{2\pi}{\omega}$。}
\bkword{単振動}{位置が $A\sin\omega t$ や $A\cos\omega t$ で表されるゆれ。$\ddot{x}=-\omega^2x$ をみたす。}
\bkword{振幅・角振動数}{ゆれの大きさ $A$、1秒に進む角 $\omega$。}
\bkword{実効値}{交流と同じ平均の電力になる直流の値。山の高さの $\frac{1}{\sqrt2}$ 倍。}
\end{kotoba}

%% ---------------- 例題 ----------------
\bksection{reidai}{例題でたしかめよう}

\begin{reidai}[$\sin$・$\cos$ をふくむ関数の微分]
次の関数を微分しましょう。\quad（1）$3\sin2t$\quad（2）$-\cos5t$\quad（3）$\sin(3t+1)$\quad（4）$\sin t^2$\par
\kangaekata 内側（角の部分）を $u$ とし、「外側の微分×内側の微分」（5-6）。前の数（定数倍）はそのまま残す（5-4）。\par
\shiki（1）内側 $u=2t$、内側の微分は2。$3\times\cos2t\times2=6\cos2t$\par
\hspace*{2.6\zw}（2）内側 $u=5t$、内側の微分は5。$-1\times(-\sin5t)\times5=5\sin5t$\par
\hspace*{2.6\zw}（3）内側 $u=3t+1$、内側の微分は $(3t+1)'=3$。$\cos(3t+1)\times3$\par
\hspace*{2.6\zw}（4）内側 $u=t^2$、内側の微分は $2t$。$\cos t^2\times2t$\par
\kotae （1）$6\cos2t$\quad（2）$5\sin5t$\quad（3）$3\cos(3t+1)$\quad（4）$2t\cos t^2$
\end{reidai}
（2）では、$\cos$ の微分の $-$ と、はじめからある $-$ で、$-$ が2つになって $+$ になります。（3）の $+1$ は定数なので、微分すると0。内側の微分は3だけです。

\begin{jibun}
次の関数を微分しましょう。\quad（1）$4\cos3t$\quad（2）$\sin\dfrac{t}{2}$\quad（3）$\cos(2t+1)$\quad（4）$-2\sin4t$
\bkjans{（1）\nobreak $-12\sin3t$\quad（2）\nobreak $\dfrac{1}{2}\cos\dfrac{t}{2}$\quad（3）\nobreak $-2\sin(2t+1)$\quad（4）\nobreak $-8\cos4t$}
\end{jibun}

\begin{reidai}[ばねのおもりの速度と加速度]
ばねのおもりの位置が $x=0.1\sin4t$（m）です。速度と加速度を求め、速さが最大になる位置とその速さ、加速度の大きさが最大になる位置とその大きさを答えましょう。\par
\kangaekata 速度は位置の微分、加速度は速度の微分（5-5）。$\sin$ も $\cos$ も、大きさの最大は1。\par
\shiki $v=0.1\times4\cos4t=0.4\cos4t$\par
\hspace*{2.6\zw}$a=0.4\times(-4\sin4t)=-1.6\sin4t$\par
\hspace*{2.6\zw}速さ：$\cos4t=\pm1$ のとき最大の0.4。このとき $\sin4t=0$ で $x=0$\par
\hspace*{2.6\zw}加速度：$\sin4t=\pm1$ のとき最大の1.6。このとき $x=\pm0.1$\par
\kotae 中心（$x=0$）を通るとき速さ最大\qty{0.4}{m/s}。はし（$x=\pm\qty{0.1}{m}$）で加速度の大きさ最大\qty{1.6}{m/s^2}
\end{reidai}
$\cos4t=\pm1$ のとき $\sin4t=0$ になるのは、$\sin^2+\cos^2=1$ だからです。速さが最大の所と、加速度が最大の所は、⑥の図のように $\frac{1}{4}$ 周期ずれています。

\begin{jibun}
$x=0.05\cos10t$（m）の速度・加速度と、速さの最大値・加速度の大きさの最大値を求めましょう。
\bkjans{$v=-0.5\sin10t$、$a=-5\cos10t$。速さ最大\qty{0.5}{m/s}、加速度最大\qty{5}{m/s^2}}
\end{jibun}

\begin{reidai}[円運動の速さと加速度]
半径\qty{0.2}{m}の円の上を毎秒5回転する点の、角速度 $\omega$ と速さを求めましょう。$x$ 座標 $x=0.2\cos\omega t$ の加速度の大きさの最大値も求めましょう。\par
\kangaekata 1回転は $2\pi$\,rad（4-2）。速さは $r\omega$（⑤）。加速度は、$x$ を2回微分して求めます（⑥）。\par
\shiki 1秒に5回転なので、1秒に回る角は $\omega=5\times2\pi=10\pi$\par
\hspace*{2.6\zw}速さは $r\omega=0.2\times10\pi=2\pi\approx6.3$\par
\hspace*{2.6\zw}$\dot{x}=-0.2\omega\sin\omega t$、$\ddot{x}=-0.2\omega^2\cos\omega t$\par
\hspace*{2.6\zw}$\cos\omega t=\pm1$ のとき最大で、$0.2\times(10\pi)^2=0.2\times100\pi^2=20\pi^2$\par
\kotae $\omega=10\pi\approx31$\,rad/s、速さ約\qty{6.3}{m/s}、加速度最大約\qty{197}{m/s^2}
\end{reidai}

\begin{jibun}
半径\qty{0.3}{m}の円の上を、1秒に1回転する点の $\omega$ と速さ、$x=0.3\cos\omega t$ の加速度の大きさの最大値を求めましょう（$\pi^2\approx9.87$）。
\bkjans{$\omega=2\pi$\,rad/s、速さ $0.6\pi\approx\qty{1.9}{m/s}$、加速度最大 $0.3\times4\pi^2\approx\qty{11.8}{m/s^2}$}
\end{jibun}

\begin{reidai}[$\sin$・$\cos$ の定積分]
$\displaystyle\int_0^{\frac{\pi}{2}}\cos t\,\dd t$ と $\displaystyle\int_0^{\pi}\sin2t\,\dd t$ を求めましょう。\par
\kangaekata ⑦の積分の式で原始関数を作り、上端の値から下端の値をひきます（5-11）。$\sin2t$ は $\omega=2$ なので $\frac{1}{2}$ がかかる。\par
\shiki $\displaystyle\int_0^{\frac{\pi}{2}}\cos t\,\dd t=\Bigl[\,\sin t\,\Bigr]_0^{\frac{\pi}{2}}=\sin\frac{\pi}{2}-\sin0=1-0=1$\par\vspace{2pt}
\hspace*{2.6\zw}$\displaystyle\int_0^{\pi}\sin2t\,\dd t=\Bigl[\,-\frac{1}{2}\cos2t\,\Bigr]_0^{\pi}=\Bigl(-\frac{1}{2}\cos2\pi\Bigr)-\Bigl(-\frac{1}{2}\cos0\Bigr)$\par\vspace{2pt}
\hspace*{4.6\zw}$\displaystyle=-\frac{1}{2}+\frac{1}{2}=0$\par
\kotae 1、0
\end{reidai}
2つめが0になるのは、$\sin2t$ が0から $\pi$ でちょうど1周期で、横軸より上の山（$+1$）と下の谷（$-1$）の面積が打ち消しあうからです（5-11）。

\begin{jibun}
$\displaystyle\int_0^{\pi}\sin t\,\dd t$ と $\displaystyle\int_0^{\frac{\pi}{4}}\cos2t\,\dd t$ を求めましょう。
\bkjans{2、$\Bigl[\,\dfrac{1}{2}\sin2t\,\Bigr]_0^{\frac{\pi}{4}}=\dfrac{1}{2}\sin\dfrac{\pi}{2}-0=\dfrac{1}{2}$}
\end{jibun}

\begin{reidai}[$\ddot{x}=-\omega^2x$ を確かめる]
$x=3\cos2t$ が、$\ddot{x}=-4x$ をみたすことを確かめましょう。\par
\kangaekata $x$ を2回微分した $\ddot{x}$ が、$-4x$ と同じ式になるかを見ます（⑥）。\par
\shiki $\dot{x}=3\times(-\sin2t)\times2=-6\sin2t$\par
\hspace*{2.6\zw}$\ddot{x}=-6\times\cos2t\times2=-12\cos2t$\par
\hspace*{2.6\zw}$-12\cos2t=-4\times3\cos2t=-4x$\par
\kotae みたす（$\omega=2$、$\omega^2=4$）
\end{reidai}

\begin{jibun}
$x=5\sin3t$ は $\ddot{x}=-kx$ をみたします。$k$ の値を求めましょう。
\bkjans{$\dot{x}=15\cos3t$、$\ddot{x}=-45\sin3t=-9x$ より $k=9$}
\end{jibun}

\begin{reidai}[観覧車の高さの変わる速さ]
半径\qty{50}{m}、1周\qty{20}{min}の観覧車があります。ゴンドラの高さは、いちばん下を出発してから $t$ 秒後に $y=50-50\cos\omega t$（m）です。上がる速さの最大値と、それが何分後かを求めましょう。最も速く下がるのは何分後ですか。\par
\kangaekata 1周の時間を秒に直して $\omega$ を求める。上がる速さは、高さ $y$ を時間で微分した $\dot{y}$（5-5）。\par
\shiki 1周は $\qty{20}{min}\times\qty{60}{s/min}=\qty{1200}{s}$\par
\hspace*{2.6\zw}1200秒で $2\pi$ 回るので $\omega=\dfrac{2\pi}{1200}$\par
\hspace*{2.6\zw}$\dot{y}=0-50\times(-\omega\sin\omega t)=50\omega\sin\omega t$\par
\hspace*{2.6\zw}最大は $\sin\omega t=1$ のとき\par
\hspace*{4.6\zw}$50\omega=50\times\dfrac{2\pi}{1200}=\dfrac{\pi}{12}\approx0.26$\par
\hspace*{2.6\zw}$\omega t=\dfrac{\pi}{2}$ は $2\pi$ の $\dfrac{1}{4}$ なので、$\qty{20}{min}\times\dfrac{1}{4}=\qty{5}{min}$\par
\hspace*{2.6\zw}最も速く下がるのは $\sin\omega t=-1$ のとき\par
\hspace*{4.6\zw}$\omega t=\dfrac{3\pi}{2}$（$\dfrac{3}{4}$ 周）で、$\qty{20}{min}\times\dfrac{3}{4}=\qty{15}{min}$\par
\kotae 上がる速さの最大は約\qty{0.26}{m/s}で5分後（真横を通るとき）。最も速く下がるのは15分後
\end{reidai}

\begin{jibun}
半径\qty{30}{m}、1周\qty{15}{min}の観覧車で、上がる速さの最大値を求めましょう。また、最も速く下がるのは何分何秒後ですか。
\bkjans{$\qty{15}{min}\times\qty{60}{s/min}=\qty{900}{s}$。$30\times\dfrac{2\pi}{900}=\dfrac{\pi}{15}\approx\qty{0.21}{m/s}$。最も速く下がるのは $15\times\dfrac{3}{4}=11.25$ 分後で、11分15秒後}
\end{jibun}

\begin{reidai}[$\sin^2$・$\cos^2$ の積分]
$\displaystyle\int_0^{\pi}\sin^2t\,\dd t$ を求めましょう。\par
\kangaekata 2倍角の式で、$\sin^2t$ を積分できる $\cos2t$ の形に直す（⑨）。\par
\shiki $\sin^2t=\dfrac{1-\cos2t}{2}=\dfrac{1}{2}-\dfrac{1}{2}\cos2t$\par\vspace{2pt}
\hspace*{2.6\zw}$\displaystyle\int_0^{\pi}\Bigl(\frac{1}{2}-\frac{1}{2}\cos2t\Bigr)\dd t=\Bigl[\,\frac{t}{2}-\frac{\sin2t}{4}\,\Bigr]_0^{\pi}=\Bigl(\frac{\pi}{2}-\frac{\sin2\pi}{4}\Bigr)-0=\frac{\pi}{2}$\par
\kotae $\dfrac{\pi}{2}$（区間の長さ $\pi$ の $\frac{1}{2}$。平均が $\frac{1}{2}$ だから）
\end{reidai}

\begin{jibun}
$\displaystyle\int_0^{\frac{\pi}{2}}\cos^2t\,\dd t$ を求めましょう。
\bkjans{$\displaystyle\int_0^{\frac{\pi}{2}}\frac{1+\cos2t}{2}\,\dd t=\Bigl[\,\frac{t}{2}+\frac{\sin2t}{4}\,\Bigr]_0^{\frac{\pi}{2}}=\frac{\pi}{4}$}
\end{jibun}

\begin{reidai}[交流の実効値と電力]
（1）実効値\qty{100}{V}のコンセントに、\qty{20}{\ohm}の抵抗をつなぎます。電圧の最大値 $V_0$ と、平均の電力を求めましょう。\par
（2）電圧の最大値が\qty{200}{V}の交流の実効値を求めましょう（$\sqrt2=1.41$）。\par
\kangaekata ⑨より、$V_0=\sqrt2\,V_{\mathrm{e}}$、$V_{\mathrm{e}}=\dfrac{V_0}{\sqrt2}$。平均の電力は、実効値と同じ電圧の直流のときと同じ。\par
\shiki（1）$V_0=\sqrt2\times100=1.414\times100\approx141$\par
\hspace*{4.6\zw}平均の電力は $\dfrac{V_{\mathrm{e}}^2}{R}=\dfrac{100^2}{20}=\dfrac{10000}{20}=500$\par
\hspace*{2.6\zw}（2）$V_{\mathrm{e}}=\dfrac{200}{\sqrt2}=\dfrac{200}{1.41}\approx141.8$\par
\kotae （1）約\qty{141}{V}、\qty{500}{W}\quad（2）約\qty{142}{V}
\end{reidai}

\begin{jibun}
最大値\qty{170}{V}の交流の実効値を求めましょう。また、\qty{50}{\ohm}の抵抗での平均の電力を求めましょう（$\sqrt2=1.41$）。
\bkjans{$V_{\mathrm{e}}=\dfrac{170}{1.41}\approx\qty{120}{V}$、平均の電力 $\dfrac{V_0^2}{2R}=\dfrac{170^2}{2\times50}=\qty{289}{W}$}
\end{jibun}

%% ---------------- まちがえやすいところ・物理工学 ----------------
\bksection{tsumazuki}{まちがえやすいところをチェック}

\begin{tsumazuki}{$\cos$ の微分の $-$ をわすれる}
\bkng{$(\cos t)'=\sin t$}\bkok{$(\cos t)'=-\sin t$}
$t=0$ の少しあとでは、$\cos t$ は1からへっていきます（傾きが負）。$\sin t$ は正なので、$-$ がないと合いません。グラフを思いうかべて確かめましょう。
\end{tsumazuki}

\begin{tsumazuki}{$\omega$ を前に出しわすれる・積分でかけてしまう}
\bkng{$\displaystyle\int\cos3t\,\dd t=3\sin3t+C$}\bkok{$\displaystyle\int\cos3t\,\dd t=\frac{1}{3}\sin3t+C$}
微分では $\omega$ をかけ、積分では $\omega$ でわります。答えを微分して確かめましょう。$\Bigl(\dfrac{1}{3}\sin3t\Bigr)'=\dfrac{1}{3}\times3\cos3t=\cos3t$ で、もとにもどります。
\end{tsumazuki}

\begin{tsumazuki}{$\sin^2t$ と $\sin t^2$ を取りちがえる}
$\sin^2t$ は $(\sin t)^2$（$\sin t$ を2乗する）。$\sin t^2$ は $\sin(t^2)$（角が $t^2$）。$t=2$ なら、$\sin^2t=(\sin2)^2=0.83$、$\sin t^2=\sin4=-0.76$ で、まったくちがう数です。
\end{tsumazuki}

\begin{tsumazuki}{電卓や表計算を「度」のまま使う}
$\sin1$ は、度のモードで0.0175、ラジアンのモードで0.8415。式に入れるのはラジアンの値です。電卓の表示（DEG／RAD）を確かめましょう。
\end{tsumazuki}

\begin{butsuri}
\begin{itemize}[leftmargin=1.2\zw, itemsep=1pt, topsep=0pt, label={\color{bkPurple}●}]
\item \textbf{ばねと振り子}\par\noindent 速さは中心で最大、はしで0。加速度ははしで最大、中心で0。$\sin$ と $\cos$ の $\frac{1}{4}$ 周期のずれが、このちがいを生みます。
\item \textbf{交流とコンデンサー}\par\noindent 家のコンセントの電圧は $V_0\sin\omega t$ の交流です。コンデンサーに流れる電流は、電圧の変わる速さ（電圧の微分）に比例するので、$\cos\omega t$ の形になり、電圧より $\frac{1}{4}$ 周期先に進みます（位相のずれ、6-10）。\par
\begin{tikzpicture}[tl/.style={font=\sffamily\footnotesize, text=bkInk, anchor=base, yshift=-6mm}]
\begin{axis}[bkplot, width=118mm, height=31mm, xmin=0, xmax=6.8, ymin=-1.15, ymax=1.15, xtick={1.5708,3.1416,4.7124,6.2832}, xticklabels={}, ytick=\empty, xlabel={$\omega t$}]
  \addplot[bkcurve, domain=0:6.6, samples=120] {sin(deg(x))};
  \addplot[bkcurve2, domain=0:6.6, samples=120] {0.8*cos(deg(x))};
  \node[tl] at (axis cs:1.5708,-1.15) {$\dfrac{\pi}{2}$};
  \node[tl] at (axis cs:3.1416,-1.15) {$\pi$};
  \node[tl] at (axis cs:4.7124,-1.15) {$\dfrac{3\pi}{2}$};
  \node[tl] at (axis cs:6.2832,-1.15) {$2\pi$};
  % 「電圧」「電流」は、格子の線が通らないように、格子の上（ymax の外）に置く
  \node[font=\sffamily\scriptsize, text=bkBlue, anchor=south west, inner sep=1pt, yshift=1.5mm] at (axis cs:1.70,1.15) {電圧};
  \node[font=\sffamily\scriptsize, text=bkOrange!80!black, anchor=south east, inner sep=1pt, yshift=1.5mm] at (axis cs:6.25,1.15) {電流（電圧より $\frac{1}{4}$ 周期先に山が来る）};
\end{axis}
\end{tikzpicture}
\item \textbf{音とスピーカー}\par\noindent スピーカーの振動板は $\sin$ の形でゆれて音を出します。振動板の加速度の大きさは $A\omega^2$ で、振動数の2乗に比例するので、同じ振幅なら、高い音ほど大きな加速度が必要です。
\end{itemize}
\end{butsuri}

%% ---------------- 章末問題 ----------------
\bkendproblems
\bkproblemgroup{基本}{3}

\bkproblem 次の関数を微分しましょう。\quad（1）$2\sin5t$\quad（2）$-\cos4t$\quad（3）$\sin(3t-1)$
\bkpans{（1）\nobreak $10\cos5t$\quad（2）\nobreak $4\sin4t$\quad（3）\nobreak $3\cos(3t-1)$}

\bkproblem $x=0.2\cos5t$（m）のおもりについて、速度と加速度を求めましょう。速さの最大値はいくつですか。
\bkpans{$v=-\sin5t$、$a=-5\cos5t$。速さ最大\qty{1}{m/s}}

\bkproblem 次の定積分を求めましょう。\\
（1）$\displaystyle\int_0^{\pi}\sin t\,\dd t$\qquad（2）$\displaystyle\int_0^{\frac{\pi}{2}}\cos t\,\dd t$\qquad（3）$\displaystyle\int_0^{\frac{\pi}{3}}\sin3t\,\dd t$
\bkpans{（1）\nobreak 2\quad（2）\nobreak 1\quad（3）\nobreak \mbox{$\Bigl[\,-\dfrac{1}{3}\cos3t\,\Bigr]_0^{\frac{\pi}{3}}=\dfrac{2}{3}$}}

\bkproblem $x=2\sin6t$ が $\ddot{x}=-kx$ をみたすときの $k$ を求めましょう。
\bkpans{$k=36$}

\bkproblem $\displaystyle\int_0^{2\pi}\cos^2t\,\dd t$ を求めましょう。また、$\cos^2t$ の1周期の平均はいくつですか。
\bkpans{$\Bigl[\,\dfrac{t}{2}+\dfrac{\sin2t}{4}\,\Bigr]_0^{2\pi}=\pi$、平均 $\dfrac{1}{2}$}

\bkproblemgroup{チャレンジ}{4}

\bkproblem 半径\qty{0.5}{m}の円の上を、1秒に2回転する点があります。$\omega$ と速さを求めましょう。また、点の $x$ 座標 $x=0.5\cos\omega t$ の加速度の大きさの最大値を求めましょう。
\bkpans{$\omega=4\pi$\,rad/s、速さ $0.5\times4\pi=2\pi\approx\qty{6.3}{m/s}$。加速度最大 $0.5\times16\pi^2=8\pi^2\approx\qty{79}{m/s^2}$}

\bkproblem 半径\qty{40}{m}、1周\qty{16}{min}の観覧車があります。いちばん下を出発してから $t$ 秒後のゴンドラの高さは $y=40-40\cos\omega t$（m）です。上がる速さの最大値を求めましょう。また、高さが最も速く下がるのは何分後ですか。
\bkpans{\mbox{$\qty{16}{min}\times\qty{60}{s/min}=\qty{960}{s}$}、\mbox{$\omega=\dfrac{2\pi}{960}$}。最大 $40\omega=\dfrac{\pi}{12}\approx\qty{0.26}{m/s}$。最も速く下がるのは $\frac{3}{4}$ 周の12分後}

\bkproblem 実効値\qty{200}{V}の交流の、電圧の最大値を求めましょう。また、\qty{40}{\ohm}の抵抗での平均の電力を求めましょう。
\bkpans{$200\sqrt2\approx\qty{283}{V}$、$\dfrac{200^2}{40}=\qty{1000}{W}$}

\begin{furikaeri}
\begin{checklist}
\item 円運動の速度で、$\sin$・$\cos$ の微分を説明できる。\hfill{\footnotesize\sffamily\color{bkTheme}→\pageref{bk:5-15:en}ページ}
\item $\sin\omega t$、$\cos\omega t$ を微分・積分できる。\hfill{\footnotesize\sffamily\color{bkTheme}→\pageref{bk:5-15:omega}ページ}
\item 単振動の $\ddot{x}=-\omega^2x$ を確かめられる。\hfill{\footnotesize\sffamily\color{bkTheme}→\pageref{bk:5-15:nikai}ページ}
\item $\sin^2$ を積分して、交流の実効値を求められる。\hfill{\footnotesize\sffamily\color{bkTheme}→\pageref{bk:5-15:jikko}ページ}
\end{checklist}
\end{furikaeri}

\bkfukidashi{mina}{$e$ も $\sin$ も、微分の計算ができるようになった。次は、積分のむずかしい形にちょうせんだよ。}
